\clearpage
\sscp{Issues in the morphosyntax}{07-04-03}
There are several issues in the morphosyntax that distinguish
\tsc{\ili{Ambon}-Seram} from \tsc{Sula-Buru} and convey the sense to anyone working with languages in both
subgroups that the two are very different.
\tsc{\ili{Ambon}-Seram} languages have Actor-Undergoer inflection or
indexing on active verbs
and some experiencer verbs,
multiple inflectional paradigms, a distinction between human
and non-human forms in third person bound pronouns,
and a lexical and grammatical distinction between alienable and inalienable nouns.
\tsc{Sula-Buru} languages do not appear to have these features.

The morphosyntax of both \tsc{Sula-Buru}
and \tsc{\ili{Ambon}-Seram} languages is built around mechanisms that
distinguish verbs as Active or non-Active,
and distinguish the semantic alignment of core arguments as being
in the macro-role of Actor or Undergoer, or both.
The underlying system behind languages in both subgroups seems
to have been based on the semantic opposition of a Split-S system \citep{Dixon1994}.
While modern \ili{Buru} no longer functions as a Split-S system,
there are traces of the current system having developed from
an earlier Split-S system \citep[156--159]{Grimes1991c}.
See \citet{Ewing2010}, and \citet{Grimes2024} for a fuller survey
of semantic opposition (also called ``semantic alignment'')
and the notion of grammatical Subject in central Maluku languages.\footnote{
		The Active/non-Active distinction in the verbs is pervasive in the
		Austronesian languages of eastern Indonesia, Timor-Leste and SHWNG.
		Split-S languages of various types are also found in \tsc{\ili{Bima}-Lembata},
		\tsc{Timor-Babar}, \tsc{Aru}, \tsc{Seram-Tanimbar-Bomberai},
		and SHWNG languages, as well as west of the Blust line.}

The Split-S variations in \tsc{\ili{Ambon}-Seram} languages on Seram
and \ili{Ambon}-Lease include some with a pronominal marking in which
the single argument S\sub{A} of Active intransitive verbs takes a different
pronominal set than the S\sub{U} of non-Active verbs.
In languages like \ili{Yalahatan},
S\sub{A} is marked \it{before} the verb with the Actor set of pronouns,
whereas S\sub{U} is marked \it{following} non-Active verbs with the
Undergoer set of pronouns.
Semantic alignment of this sort found in \tsc{\ili{Ambon}-Seram} languages
is relatively uncommon in the languages of the world \citep{Siewierska2013},
although it is widespread in our target region,
it is also sporadic,
and is not a “feature” or “innovation” of the region that can
be used for subgrouping purposes at a higher level.

\ili{Yalahatan} active transitive clauses have the subject-Actor \it{preceding}
the verb, and the object-Undergoer \it{following}
the verb as in examples \qf{07-05} and \qf{07-06}.

\begin{exe}
\setcounter{xnumi}{4}
	\ex{\gll	si=hita ulu-si\\
						\tsc{3pl.a}=chip head-\tsc{3pl.gen}\\}
						\jambox*{(\ili{Yalahatan}: active transitive)}
			\glt	‘They\sub{\it{i}} cut off their\sub{\it{j}} heads.’
						\label{ex:07-05}
	\ex{\gll	yaʔu waru=i\\
						\tsc{1sg} hit-\tsc{3sg.u}\\}
						\jambox*{(\ili{Yalahatan}: active transitive)}
			\glt	‘I hit him.’
						\label{ex:07-06}
\end{exe}

The single S\sub{U} argument of a non-active intransitive verb is marked
by subject indexing \it{following} the verb,
in boldface in example \qf{07-07}.

\begin{exe}
	\ex{\gll	{\ldots} lau mata=\tbf{ʔu}\\
						{} until die=\tsc{1sg.u}\\}
						\jambox*{(\ili{Yalahatan}: non-active intransitive)}
			\glt	‘{\ldots} until I die.’
						\label{ex:07-07}
\end{exe}

Some languages in the central Maluku region,
such as \ili{Sou Amana Teru} \citep{Musgrave2010}
and \ili{Yalahatan} have verbs of motion,
posture, bodily function, and sometimes experiencer verbs,
which are semantically intransitive (one referent),
but grammatically (optionally) transitive,
with Actor and Undergoer coreferential -- the one doing the action
is also the one undergoing or experiencing the action or quality
of the verb, or the one whose location is being changed,
as in example \qf{07-08}.

\begin{exe}
	\ex{\gll	amite \tbf{yaʔu} tue=\tbf{ʔu} {\ldots}\\
						last.night \tsc{1sg} sit=\tsc{1sg.u}\\}
						\jambox*{(\ili{Yalahatan}: posture)}
			\glt	‘Last night I was sitting {\ldots}’
						\label{ex:07-08}
\end{exe}

Many experiencer verbs are also syntactically transitive with
coreferential Subject and Object, as in example \qf{07-09}.
The agency of the Subject-experiencer is low.

\begin{exe}
	\ex{\gll	\tbf{yaʔu} atere=\tbf{ʔu}.\\
						\tsc{1sg} afraid=\tsc{1sg.u}\\}
						\jambox*{(\ili{Yalahatan}: experiencer)}
			\glt	‘I'm afraid.’
						\label{ex:07-09}
\end{exe}

These should not be interpreted as causative or reflexive (i.e.
‘I made myself afraid’,
‘I frightened myself’), as those would take morphological or
periphrastic causatives or explicit reflexive pronouns.
Compare \ili{Buru} examples \qf{07-10}--\qf{07-11} with \ili{Yalahatan} example \qf{07-09}.

\begin{exe}
	\ex{\gll	\tbf{yako} sei \tbf{kono} haik fi di nani fina emsawan rua naa.\\
						\tsc{1sg} weary \tsc{1sg.u} \tsc{pfv} \tsc{loc} \tsc{dist} \tsc{1pi.poss} female C-in-law two \tsc{prox}\\}
			\glt	‘I am fed up with these two daughters-in-law of ours.’ \hfill{(\ili{Buru}: experiencer)}
						\label{ex:07-10}
	\ex{\gll	\tbf{yako} la ep-toke emhewak \tbf{kono}, {\ldots}\\
						\tsc{1sg} \tsc{irr} \tsc{caus}-show self \tsc{1sg.u}\\}
						\jambox*{(\ili{Buru}: reflexive)}
			\glt	‘I want to teach myself, [how to read and write \ili{Buru}].'
						\label{ex:07-11}
\end{exe}

\citet{Klamer2008} and \citet{Schapper2015} note that Split-S
marking dividing intransitive verbs into Active
and non-Active clause types are found in both the Austronesian
and non-Austronesian languages of eastern Indonesia.
Interestingly enough verbs of motion
and posture optionally marked as transitive,
found in many Austronesian languages of eastern Indonesia
and Timor-Leste, are also found in the non-Austronesian (Papuan)
languages of the region.
So once again, contact cannot be ruled out as the source for
this type of verb marking in the Austronesian languages.
The data in \qf{07-12} are from Western \ili{Pantar} (a Papuan Timor-\ili{Alor}-\ili{Pantar}
language).

\begin{exe}
	\ex{\gll	naŋ na-lama ta.\\
						\tsc{1sg} \tsc{1sg.u}-go \tsc{ipfv}\\}
						\jambox*{(Western \ili{Pantar}: motion; \citealp[105]{Holton2010a})}
			\glt	‘I'm going.’
						\label{ex:07-12}
\end{exe}

A similar pattern of “intransitive” motion
and posture verbs with coreferential Undergoer pronouns is found
in \ili{Klon} \citep[50--51, 81--84]{Banamtuan2021} as in examples \qf{07-13}--\qf{07-14},
a Timor-\ili{Alor}-\ili{Pantar} language of \ili{Alor}.

\begin{exe}
	\ex{\gll	naan na agai.\\
						\tsc{1sg} \tsc{1sg.u} go\\}
						\jambox*{(\ili{Klon}: motion; \citealp[84]{Banamtuan2021})}
			\glt	‘I'm going now.’
						\label{ex:07-13}
	\ex{\gll	ga a moop.\\
						\tsc{3sg} \tsc{3sg.u} sleep\\}
						\jambox*{(\ili{Klon}: posture; \citealp[51]{Banamtuan2021})}
			\glt	‘He slept.’
						\label{ex:07-14}
\end{exe}

In the morphosyntax, there are three related issues that distinguish
\tsc{\ili{Ambon}-Seram} languages from \tsc{Sula-Buru} languages.

\begin{itemize}
	\item	\tsc{\ili{Ambon}-Seram} languages tend to inflect verbs for person
				and number with Subject (or Actor)
				and Object (or Undergoer) indexes on verbs, particularly active verbs.
				In \tsc{Sula-Buru}, \tsc{Sula} inflects active verbs.
				\tsc{Buru} does not inflect its verbs.
	\item	\tsc{\ili{Ambon}-Seram} languages can have several different inflectional
				paradigms (\citet{Collins1982c,Collins1983} note up to seven paradigms in single languages),
				several of which can involve morphophonemic processes that mask
				the underlying consonant of the historical root.
				\tsc{Sula-Buru} languages do not
				(see \trf{07-16}).
	\item	\tsc{\ili{Ambon}-Seram} languages distinguish between human
				and non-human in their third person forms.
				While the data are sparse,
				there is no clear evidence that \tsc{Sula-Buru} languages make
				such a distinction (see discussion and footnote below).
\end{itemize}

A fourth issue is independent of the previous three.

\begin{itemize}
	\item	\tsc{\ili{Ambon}-Seram} languages distinguish between alienable
				and inalienable nouns in their lexicon and grammars.
				\tsc{Sula-Buru} languages do not make such a distinction (see
				\citet[277-292]{Grimes1991c}: for \tsc{Buru};
				\citet[280--293]{Bloyd2020}: for \tsc{Sula}).
\end{itemize}

\ili{As} mentioned in \srf{01-07},
it is important for comparative purposes to note that the inflectional
paradigms for verbs in \tsc{\ili{Ambon}-Seram} languages are only partially
cognate with the inflectional paradigms of \tsc{Timor-Babar},
\tsc{Seram-Tanimbar-Bomberai}, or SHWNG languages.
The forms that are cognate link back to retentions of PMP forms,
so \it{contra} \citet[258]{Blust1993}, subject markers on the
verb are not an innovation supporting CEMP as a subgroup.
This is revisited in \srf{13-03-03}.

\ili{As} mentioned above, \tsc{\ili{Ambon}-Seram} languages tend to inflect
verbs for person and number \citep[24--26]{Collins1982c,Collins1983}.
And each language can have several different inflectional paradigms
that involve different morphophonemic alternations.
\citet[118--125]{Stresemann1927} provides inflectional paradigms for
several verbs for a number of \tsc{\ili{Ambon}-Seram} languages.
In \ili{Asilulu} \citep[xxix--xxx]{Collins2003a} these are restricted to Active verbs.
\citet[36--42]{Bolton1990} describes the verb inflections for South \ili{Nuaulu}.
\citet[86--91]{Laidig1992} describes these for \ili{Larike},
with some morphophonemic alternations of the initial {\#}C of
the verb root illustrated in \trf{07-16}.

\begin{table}\tcp{Subject indexing in \ili{Larike} showing morphophonemics}{07-16}
	\begin{threeparttable}\st{6pt}
	\begin{tabular}{l|r@{ }llll}\lsptoprule
paradigm				&				&	p-verb			&	k-verb			&	s-verb			&	t-verb	\\	
gloss						&				&	work				&	bite				&	ascend			&	know	\\	
PMP							&				&							&	*kitkit			&	*sakay			&	**tewa\su{†}	\\	\midrule
\tsc{1sg}				&	au-		&	-\tbr{ʔ}ese	&	-\tbr{ʔ}iʔi	&	-\tbb{s}aʔa	&	\hr-\tbb{t}iwa	\\	
\tsc{2sg}				&	ai-		&	-\tbb{p}ese	&	-\tbb{k}iʔi	&	-\tbr{r}aʔa	&	\hr-\tbr{r}iwa	\\	
\tsc{3sg.hum}		&	mei-	&	-\tbb{p}ese	&	-\tbb{k}iʔi	&	-\tbr{r}aʔa	&	\hr-\tbr{r}iwa	\\	
\tsc{3sg.nhum}	&	i-		&	-\tbb{p}ese	&	-\tbb{k}iʔi	&	-\tbr{r}aʔa	&	\hr-\tbr{r}iwa	\\	
\tsc{1pl.incl}	&	ite-	&	-\tbr{ʔ}ese	&	-\tbr{ʔ}iʔi	&	-\tbb{s}aʔa	&	\hr-\tbb{t}iwa	\\	
\tsc{1pl.excl}	&	ami-	&	-\tbr{ʔ}ese	&	-\tbr{ʔ}iʔi	&	-\tbb{s}aʔa	&	\hr-\tbb{t}iwa	\\	
\tsc{2pl}				&	imi-	&	-\tbr{ʔ}ese	&	-\tbr{ʔ}iʔi	&	-\tbb{s}aʔa	&	\hr-\tbb{t}iwa	\\	
\tsc{3pl.hum}		&	mati-	&	-\tbr{ʔ}ese	&	-\tbr{ʔ}iʔi	&	-\tbb{s}aʔa	&	\hr-\tbb{t}iwa	\\	
\tsc{3pl.nhum}	&	iri-	&	-\tbb{p}ese	&	-\tbb{k}iʔi	&	-\tbr{r}aʔa	&	\hr-\tbr{r}iwa	\\	
		\lspbottomrule
	\end{tabular}
		\begin{tablenotes}\tnsize
			\item[†] {Compare \ili{Buru} \it{tewa} ‘know, understand’,
								\ili{Asilulu} \it{tewa} ‘know’, \ili{Alune} \it{tekwa} ‘know’.}
		\end{tablenotes}
	\end{threeparttable}
\end{table}

\newpage
Apart from \tsc{\ili{Ambon}-Seram} languages,
similar morphophonemic changes in the initial consonants also
occur in the \tsc{Eastern Islands} language \ili{Geser}-\ili{Gorom},
usually involving alternations between voiceless
and voiced plosives \citep[55--57]{Stokhof1982b}.
Likewise, \citet[527--533]{CollinsKaartinen1998} describe similar
morphophonemic changes in the initial consonants of certain verbs
in \ili{Banda}, as illustrated in \trf{07-17}.
In \ili{Banda} these alternations often involve a nasal or prenasalised
consonant, though other consonants are also involved.

\begin{table}\tcp{Subject indexing in \ili{Banda} showing morphophonemics}{07-17}
	\begin{threeparttable}\st{6pt}
	\st{6pt}\begin{tabular}{lllll}\lsptoprule
paradigm	&	s-verb	&	f-verb	&	k-verb	&	t-verb	\\
gloss	&	bathe	&	pay	&	bite	&	fly	\\
PMP	&		&	*bayaD	&	*kitkit	&	**ntibu\su{†}	\\ \midrule
\tsc{1sg}	&	\tbf{ʧ}eli	&	\tbf{k}iar	&	\tbb{k}iki	&	\tbb{t}ifu	\\
\tsc{2sg}	&	\tbr{ɲʤ}eli	&	\tbr{mb}iar	&	\tbr{m(b)}iki	&	\tbr{nd}ifu	\\
\tsc{3sg}	&	\tbr{ɲʤ}eli	&	\tbr{mb}iar	&	\tbf{ŋ}iki	&	\tbr{nd}ifu	\\
\tsc{1pl.incl}	&	te-\tbb{s}eli	&	ti-\tbb{f}iar	&	ti-\tbb{k}iki	&	ti-\tbb{t}ifu	\\
\tsc{1pl.excl}	&	ke-\tbb{s}eli	&	ki-\tbb{f}iar	&	ki-\tbb{k}iki	&	ki-\tbb{t}ifu	\\
\tsc{2pl}	&	\tbr{ɲʤ}eli	&	\tbr{mb}iar	&	\tbr{m(b)}iki	&	\tbr{nd}ifu	\\
\tsc{3pl}	&	re-\tbb{s}eli	&	ri-\tbb{f}iar	&	ri-\tbb{k}iki	&	ri-\tbb{t}ifu	\\
		\lspbottomrule
	\end{tabular}
		\begin{tablenotes}\tnsize
			\item[†] {See example \qf{07-03} in \srf{07-02-02}}
		\end{tablenotes}
	\end{threeparttable}
\end{table}

In the morphosyntax, \tsc{\ili{Ambon}-Seram} languages also distinguish
\it{human} from \it{non-human} in \tsc{3sg}
and \tsc{3pl} pronominal marking,
as well as Actor (A) versus Undergoer (U),
as illustrated in \trf{07-18}.
\citet[28]{Collins1983} also provides additional data for other \tsc{\ili{Ambon}-Seram} languages.
Other languages of central Maluku do not appear to make such
a human/non-human distinction.\footnote{
		\citet{Bloyd2020} gives
		contradictory information on this for the Sula (\ili{Sanana}) language.
		He says there is a human/non-human distinction in Sula.
		But in terms of form,
		function and distribution, the “pronominal prefix” \it{i-} is
		identical for \tsc{3sg} human,
		\tsc{3sg} non-human and \tsc{3pl} (p. 233, 417, 568).
		The sentence example to illustrate the “non-human” prefix (p.234)
		places it on the verb,
		but elsewhere (p. 47) he says “verbs are not marked to agree with non-human subjects”.
		The example on p. 313 is at odds with the example on p.234.}
Thus, the human/non-human distinction sets \tsc{\ili{Ambon}-Seram}
apart from neighbouring subgroups.
Note, however, that languages of \ili{Aru} also have a human/non-human
distinction as described by \citet{Schapper2015}.
It also occurs in Papuan languages of Wallacea
and in SHWNG languages.
\citet{Schapper2015} sees “neuter gender” as one feature of
linguistic Wallacea.\footnote{
		Although \citet{Schapper2015} uses the term “neuter gender”
		descriptions of \tsc{\ili{Ambon}-Seram}
		languages use the terms human/non-human.} 

\begin{table}
	\centering\tcp{Human/non-human marking in \tsc{\ili{Ambon}-Seram} languages}{07-18}
	\st{3.2pt}\begin{tabular}{llr|lr|lr|lr|lr|lr}\lsptoprule
&\mc{2}{c|}{Larike}&\mc{2}{c|}{Asilulu}&\mc{2}{l|}{Sou Amana}&\mc{2}{c|}{Alune}&\mc{2}{c|}{S. Nuaulu}&\mc{2}{c}{Yalahatan}\\
&&&&&\mc{2}{l|}{Teru}&&&&&\\
								&	A			& U		&	A		&	U		&	A		&	U		&	A			&	U			&	A				&	U			&	A		&	U		\\ \midrule
\tsc{3sg.hum}		&	mei-	&			&	i-	&	-ni	&	e-	&	-i	&	(e)i-	&	-(n)i	&	i-			&	-(k)i	&	i-	&	-i	\\
\tsc{3sg.nhum}	&	i-		&	-a	&	a-	&	-ne	&			&	-re	&	e-		&	-(l)e	&	(e)re-	&	{\0}	&	e-	&	-e	\\
\tsc{3pl.hum}		&	mati-	&			&	si-	&	-se	&	si-	&	-si	&	si-		&	-si		&	o-			&	-so		&	si-	&	-si	\\
\tsc{3pl.nhum}	&	iri-	&			&	ru-	&	-ru	&			&			&	u-		&	-(l)u	&	(e)ra-	&	-re		&	ru-	&	-ru	\\
		\lspbottomrule
	\end{tabular}
\end{table}
